% \begin{comment}
%<*driver>
\documentclass{pildoc}

\usepackage[spanish, es-nodecimaldot]{babel}

\title{Curriculum Vitae}
\author{Joar Buitrago\thanks{jebuitragoc@unal.edu.co}}

\begin{document}
\maketitle
\part{Introducción}

Este documento está orientado a guiarte en el proceso de generar una hoja de vida profesional en {\lmr\LaTeX}.
Esta plantilla ha sido diseñada bajo el formato \texttt{.dtx} ({\lmr Documented \LaTeX}), lo que significa que el código fuente de la plantilla y este manual de usuario coexisten en un único archivo modular e independiente.
Si compilas este documento se obtiene esta guía de uso y estilo.

El objetivo principal de utilizar esta clase es separar por completo el contenido de la forma.
Mientras te concentras en redactar tus logros y experiencia de manera clara y precisa, {\lmr\LaTeX} se encarga de gestionar la consistencia visual, la jerarquía tipográfica y la optimización del espacio.
Esto garantiza un documento elegante para un evaluador humano y perfectamente legible para los sistemas automatizados de selección de personal (ATS).



\section{Estructura del proyecto y flujo de trabajo}

Esta plantilla se compone estrictamente de dos archivos fundamentales que trabajan en conjunto:

\begin{itemize}
    \item \textbf{Archivo de instalación (\texttt{.ins}):} Es el encargado de generar el motor de diseño. Al compilarlo, se extrae automáticamente el archivo de clase (\texttt{.cls}), el cual contiene toda la lógica, macros y herramientas de formato de la plantilla.
    \item \textbf{Archivo documentado (\texttt{.dtx}):} Es el archivo que estás leyendo en este momento. Su función es doble: por un lado, contiene el código fuente de la plantilla y, por el otro, genera este manual de usuario para explicarte cómo utilizar cada comando.
\end{itemize}



\section{Instrucciones para Empezar}

Para comenzar a construir tu hoja de vida, debes seguir este orden cronológico en tu entorno de desarrollo o terminal:

\begin{enumerate}
    \item \textbf{Generar la documentación:} Para generar el manual de documentación de la clase debes ejecutar el comando:
    \begin{verbatim}
lualatex cv.dtx
    \end{verbatim}

    \item \textbf{Generar la clase:} Compila el archivo \texttt{.ins} utilizando:
    \begin{verbatim}
lualatex cv.ins
    \end{verbatim}
    Esto crea el archivo \texttt{cv.cls} en tu directorio de trabajo.

    \item \textbf{Descargar las fuentes:} Antes de poder compilar el documento, es necesario descargar las fuentes que se utilizan:
    \begin{itemize}
        \item En Windows, ejecute el archivo \texttt{fonts.ps1}. Muy probablemente Windows te pida permisos para ejecutarlo.
        \item En Linux o MacOS, ejecute el archivo \texttt{fonts.sh}, para poder ejecutarlo necesitará \texttt{curl} y \texttt{jq}.
    \end{itemize}

    \item \textbf{Escribir la Hoja de Vida:} Crea un archivo nuevo con extensión \texttt{.tex} (por ejemplo, \texttt{cv.tex}) en la misma carpeta donde se encuentra el archivo \texttt{cv.cls} generado. En la primera línea del documento, invoca la clase utilizando:
    \begin{verbatim}
\documentclass{cv}
    \end{verbatim}
\end{enumerate}

A partir de este punto, puedes utilizar las siguientes partes de este manual para conocer los comandos específicos que esta plantilla te ofrece para dar formato a tus datos personales, educación, experiencia laboral y habilidades.



\section{Cómo leer esta guía}

Este manual está organizado de forma secuencial, siguiendo el mismo orden lógico en el que construyes una hoja de vida.
Comenzando por la configuración global, el encabezado con tus datos personales, y avanzando hacia las secciones.

A lo largo de las secciones, encontrarás bloques técnicos donde se definen los nuevos comandos utilizando la nomenclatura estándar de la documentación de {\lmr\LaTeX}:
\begin{itemize}
    \item Los comandos nuevos aparecen resaltados en el margen izquierdo para su rápida localización.
    \item La sintaxis te indicará los argumentos requeridos entre llaves \marg{argumento} y los opcionales (si los hay) entre corchetes \oarg{argumento}.
\end{itemize}

Te recomiendo mantener este manual abierto como referencia mientras redactas tu documento principal \texttt{.tex}, avanzando sección por sección.

\DocInput{cv.dtx}
\end{document}
%</driver>
% \end{comment}
%
%
%
%
% \begin{comment}
%<*class>
\NeedsTeXFormat{LaTeX2e}
\ProvidesClass{cv}[Curriculum Vitae LaTeX document class]



\newif\if@noicons
\DeclareOption{noicons}{\@noiconstrue}
\DeclareOption{titlepage}{\ClassError{cv}{%
   The option `titlepage' is disabled for this document class.
}{%
    You must delete it immediately from the options in your document.
}}
\DeclareOption*{\PassOptionsToClass{\CurrentOption}{article}}
\ProcessOptions\relax
\LoadClass{article}



\RequirePackage{hyperref}
\RequirePackage{hyperxmp}
\hypersetup{
    pdfusetitle,
    hidelinks,
    unicode,
}



\RequirePackage{geometry}
\geometry{margin=0.5in}
\setlength{\parindent}{0pt}
\setlength{\parskip}{0pt}



\RequirePackage[no-math]{fontspec}
\setmainfont{NotoSans}[
    Path = fonts/,
    Extension = .ttf,
    UprightFont = *-Regular,
    BoldFont = *-Bold,
    ItalicFont = *-Italic,
    BoldItalicFont = *-BoldItalic,
]
\newfontfamily{\tablerfamily}{TablerIcons}[
    Renderer = HarfBuzz,
    UprightFont = *-Outlined,
    BoldFont = *-Filled,
    Path = fonts/,
    Extension = .ttf,
]



\RequirePackage{graphicx}
\RequirePackage{tikz}



\RequirePackage{enumitem}
\setlist[itemize]{
    nosep,
    noitemsep,
    leftmargin=*,
}
%</class>
% \end{comment}
%
%
%
% \part{Creación del documento}
%
% El objetivo de esta sección es introducir la interfaz de usuario de la clase. En lugar de lidiar con entornos complejos, espaciados manuales o diseños de tablas engorrosos, esta plantilla abstrae la complejidad en comandos intuitivos y semánticos diseñados específicamente para una hoja de vida.
%
%
%
% \section{Opciones de clase}
%
% Esta plantilla o clase está basada en la clase \texttt{article} estándar de {\lmr\LaTeX}, por lo que todas las opciones habilitadas para la clase \texttt{article} estarán disponibles también en esta.
%
% Las opciones disponibles para la clase, y que te pueden ser más útiles son:
% \begin{description}
%     \item[\texttt{@pt}] Configura el tamaño de la fuente por defecto en {\lmr\LaTeX}. Los valores válidos para \texttt{@} son \texttt{10}, \texttt{11} y \texttt{12}. Por defecto está configurado en \texttt{10pt}.
%
%     \item[\texttt{@paper}] Configura el tamaño de la hoja. Los valores válidos para \texttt{@} son
%     \begin{description}
%         \item[\texttt{a4}] Alto: \qty{297}{\milli\metre} -- Ancho: \qty{210}{\milli\metre}
%         \item[\texttt{a5}] Alto: \qty{210}{\milli\metre} -- Ancho:\qty{148}{\milli\metre}
%         \item[\texttt{b5}] Alto: \qty{250}{\milli\metre} -- Ancho: \qty{176}{\milli\metre}
%         \item[\texttt{letter}] Alto: \qty{11}{\inch} -- Ancho: \qty{8.5}{\inch}
%         \item[\texttt{legal}] Alto: \qty{14}{\inch} -- Ancho: \qty{8.5}{\inch}
%         \item[\texttt{executive}] Alto: \qty{10.5}{\inch} -- Ancho: \qty{7.25}{\inch}
%     \end{description}
%     Por defecto suele estar configurado en \texttt{letterpaper}, depende de la configuración de tu distribución de {\lmr\LaTeX}.
%
%     \item[\texttt{landscape}] Configura la orientación de la hoja como horizontal.
%
%     \item[\texttt{draft}] Configura el documento en modo borrador. Muy útil para previsualizaciones rápidas.
%
%     \item[\texttt{noicons}] Deshabilita el uso de iconos en el documento.
% \end{description}
%
%
%
% \section{Título}
%
% \begin{macro}{\name}
% Empezaremos registrando tu nombre, para eso utiliza el comando \cmd\name\marg{nombre}\marg{apellido}.
% Este registra tu nombre en el documento y lo almacena para su uso futuro.
% Se recibe tu nombre y apellido como dos argumentos distintos para darles formato a cada uno de manera independiente.
% \begin{comment}
%<*class>
\NewDocumentCommand{\name}{ m m }{%
    \DeclareExpandableDocumentCommand{\@givenname}{}{#1}%
    \DeclareExpandableDocumentCommand{\@surname}{}{#2}%
}
%</class>
% \end{comment}
% \end{macro}
%
%
%
% \begin{macro}{\photo}
% Para añadir una foto junto al nombre, solamente debes llamar el comando \cmd\photo\marg{ruta}, cuyo único argumento, es la ruta de la imagen.
% La foto se incluirá utilizando el paquete \texttt{graphicx}, por lo que solo los formatos recibidos por este paquete podrán ser utilizados.
% \begin{comment}
%<*class>
\NewDocumentCommand{\photo}{ m }{\DeclareExpandableDocumentCommand{\@photo}{}{#1}}
%</class>
% \end{comment}
% \end{macro}
%
%
%
% \begin{macro}{\tagline}
% Para incluir el cargo al que aspiras, o una etiqueta, puedes utilizar el comando \cmd\tagline\marg{texto}.
% Idealmente, este debe ser un texto corto de no más de 30 caracteres.
% \begin{comment}
%<*class>
\NewDocumentCommand{\tagline}{ m }{\DeclareExpandableDocumentCommand{\@tagline}{}{#1}}
%</class>
% \end{comment}
% \end{macro}
%
%
%
% \begin{comment}
%<*class>
\let\@personalinfo\@empty
\newif\if@firstpersonalinfo \@firstpersonalinfotrue
\NewDocumentCommand{\@infosep}{}{\if@firstpersonalinfo\@firstpersonalinfofalse\else\qquad\fi}
%</class>
% \end{comment}
%
%
%
% \begin{macro}{\mail}
% Luego de esto, muy probablemente desees incluir tu correo, esto se hace con el comando \cmd\mail\marg{correo}. Este comando añade un icono, a menos de que se defina \texttt{noicons} en las opciones de la clase. Además, el contenido será interactivo, y al hacer clic sobre él, automáticamente abre el gestor de correos para enviar un mensaje a tu dirección.
% \begin{comment}
%<*class>
\NewDocumentCommand{\mail}{ m }{\DeclareExpandableDocumentCommand{\@mail}{}{#1}\personalinfo{\symbol{"EAE5}}{\href{mailto:#1}{#1}}}
%</class>
% \end{comment}
% \end{macro}
%
%
%
% \begin{macro}{\address}
% También es posible que desees incluir tu dirección o domicilio, esto lo puedes lograr mediante el comando \cmd\address\marg{dirección}. Este comando añade un icono, que de igual forma, si se define \texttt{noicons} en las opciones de la clase, desaparecerá.
% \begin{comment}
%<*class>
\NewDocumentCommand{\address}{ m }{\personalinfo{\symbol{"EAE8}}{#1}}
%</class>
% \end{comment}
% \end{macro}
%
%
%
% \begin{macro}{\phone}
% De igual forma, puedes incluir tu número telefónico con el comando \cmd\phone\marg{teléfono}. Este comando añade un icono, que nuevamente será omitido si se define \texttt{noicons} en las opciones de la clase. Además, el contenido será interactivo, y al hacer clic sobre él, automáticamente inicia el aplicativo de teléfono.
% \begin{comment}
%<*class>
\NewDocumentCommand{\phone}{ m }{\personalinfo{\symbol{"EB09}}{\href{tel:#1}{#1}}}
%</class>
% \end{comment}
% \end{macro}
%
%
%
% \DescribeMacro{\homepage}
% \DescribeMacro{\linkedin}
% \DescribeMacro{\github}
% \DescribeMacro{\gitlab}
% Existen muchísimos más comandos destinados a la inclusión de información personal, todos tienen la misma sintaxis, y todos incluyen un icono personalizado.
% Por ejemplo, el comando \cmd\homepage\marg{url} se utiliza para generar una entrada con tu dirección web para mostrar tu portafolio en línea, o un blog, o cualquier otro elemento que se pueda encontrar en internet.
% \begin{comment}
%<*class>
\NewDocumentCommand{\homepage}{ m }{\personalinfo{\symbol{"F38F}}{\href{https://#1}{#1}}}
\NewDocumentCommand{\linkedin}{ m }{\personalinfo{\symbol{"EC8C}}{\href{https://linkedin.com/in/#1}{#1}}}
\NewDocumentCommand{\github}{ m }{\personalinfo{\symbol{"EC1C}}{\href{https://github.com/#1}{#1}}}
\NewDocumentCommand{\gitlab}{ m }{\personalinfo{\symbol{"EC1D}}{\href{https://gitlab.com/#1}{#1}}}
%</class>
% \end{comment}
%
%
%
% \begin{macro}{\profileinfo}
% Pero, «¿y que hago si lo que quiero añadir no existe como comando?» En este caso, he incluido un comando especial, el comando \cmd\profileinfo\marg{icono}\marg{texto}, que te permite escoger el \textit{icono} que quieres que acompañe al \textit{texto}.
% Este comando, al igual que los demás, omite el icono en caso de que la opción de clase \texttt{noicons} haya sido utilizada en el documento.
%
% En este momento, los iconos deben escribirse como un \href{https://unicode.org/glossary/#code_point}{\emph{code point de Unicode}}, esto se logra mediante el comando \cmd\symbol\marg{code point}, en donde \meta{code point} debe escribirse en mayúsculas iniciando con comillas dobles (\texttt{"} -- hex: \texttt{22}), de manera que para escribir, por ejemplo, un \texttt{*} debes utilizar |\symbol{"2A}|.
% La fuente de iconos que utilizo en esta plantilla se llama \emph{Tabler}, que posee una colección de más de 6000 iconos distintos para que utilices en cualquier parte de tu documento, para más información revisa la sección \ref{sec:icons}.
% Puedes consultar cuáles son los iconos disponibles, así como sus \emph{code points} en la \href{https://tabler.io/icons}{\emph{página de la fuente}}.
% \begin{comment}
%<*class>
\NewDocumentCommand{\personalinfo}{ m m }{\appto\@personalinfo{\@infosep\mbox{\tableroutlined{#1}\space#2}}}
%</class>
% \end{comment}
% \end{macro}
%
%
%
% \begin{macro}{\profile}
% Finalmente, puedes incluir una descripción de tu perfil, el comando \cmd\profile\marg{texto} nos permite crear esta descripción.
% Este es un comando largo, eso quiere decir que su argumento recibir párrafos o saltos de línea.
% \begin{comment}
%<*class>
\NewDocumentCommand{\profile}{ +m }{\DeclareExpandableDocumentCommand{\@profile}{}{#1}}
%</class>
% \end{comment}
% \end{macro}
%
%
%
% \begin{macro}{\prompt}
% Hoy en día, la mayoría de ATS utilizan LLM para la generación automática de un perfil, que muchas veces no es leído por un humano.
% Teniendo en cuenta esta práctica común, el comando \cmd\prompt\marg{texto} de manera oculta en el documento, de manera que sea visible para los LLM pero no para los humanos.
% El contenido de este prompt es lo primero que lee el LLM.
% \meta{texto} es un argumento \emph{verbatim} o literal, esto quiere decir que {\lmr\LaTeX} no interpreta absolutamente nada del argumento como comando, incluyendo comentarios de {\lmr\LaTeX}.
% No recomiendo el uso de esta herramienta, ya que si alguien revisa tu documento y selecciona todo el texto, será muy notable.
% \begin{comment}
%<*class>
\NewDocumentCommand{\prompt}{ +v }{\DeclareExpandableDocumentCommand{\@prompt}{}{#1}}
%</class>
% \end{comment}
% \end{macro}
%
%
%
% \begin{macro}{\maketitle}
% El siguiente paso en el proceso de creación de la hoja de vida debe ser colocar el título. Al igual que en otras clases estándar de {\lmr\LaTeX} colocaremos el título por medio del comando {\cmd\maketitle}, este comando no recibe ningún argumento, y debería ser el primer comando que aparezca en tu entorno \texttt{document}.
% Este comando espera que hayas definido tu nombre. Si no lo haz definido, arrojara un error. Tomando toda la información que suministraste con los otros comandos de esta sección.
%
% Podrás notar también, al momento de utilizar este comando, que toda tu información personal aparecerá en el orden que tú hayas especificado con los demás comandos vistos en esta sección.
% \begin{comment}
%<*class>
\newlength{\beforetaglineskip}
\setlength{\beforetaglineskip}{0\baselineskip\@plus0.5\baselineskip}
\newlength{\beforepersonalinfoskip}
\setlength{\beforepersonalinfoskip}{1\baselineskip\@plus1\baselineskip}
\newlength{\beforeprofileskip}
\setlength{\beforeprofileskip}{0.5\baselineskip\@plus0.5\baselineskip\@minus0.5\baselineskip}
\newlength{\titleskip}
\setlength{\titleskip}{2\baselineskip\@plus0.5\baselineskip\@minus0.5\baselineskip}
\newlength{\photowidth}
\setlength{\photowidth}{1in}
\newlength{\@infowidth}
\RenewDocumentCommand{\maketitle}{}{%
    \ifx\@givenname\@undefined\ClassError{cv}{Given Name is undefined}{Use the \name command to define it on your document}\fi%
    \ifx\@surname\@undefined\ClassError{cv}{Surname is undefined}{Use the \name command to define it on your document}\fi%
    \setlength{\@infowidth}{\linewidth}%
    \ifx\@prompt\@undefined\else%
        \begin{tikzpicture}[remember picture,overlay]%
            \node at (current page.north) [anchor=north,inner sep=0pt,minimum width=\paperwidth] {%
                \parbox{\linewidth}{%
                    \color{white}\fontsize{0.1pt}{0.12pt}\selectfont\@prompt%
                }};%
        \end{tikzpicture}%
    \fi%
    \ifx\@photo\@undefined\else\addtolength{\@infowidth}{-\photowidth}\fi%
    \hypersetup{
        pdftitle = {Curriculum Vitae},%
        pdfauthor = {\@givenname\space\@surname\ifx\@mail\@undefined\else\space<\@mail>\fi},%
        pdfsubject = {\@givenname\space\@surname},%
        pdfkeywords = {Curriculum Vitae,\@givenname\space\@surname},%
        pdfcreator = {\@givenname\space\@surname\space with LaTeX},%
    }%
    \begingroup%
    \begin{minipage}{\@infowidth}%
        {\fontsize{30}{36}\selectfont\@givenname\space\bfseries\@surname\par}%
        \ifx\@tagline\@undefined\else{\vspace{\beforetaglineskip}\Large\@tagline\par}\fi%
        \ifx\@personalinfo\@empty\else{\vspace{\beforepersonalinfoskip}\raggedright\@personalinfo\par}\fi%
    \end{minipage}%
    \ifx\@photo\@undefined\else\hfill%
    \begin{minipage}{\photowidth}%
        \raggedleft%
        \begin{tikzpicture}%
            \clip circle (0.5\linewidth);
            \node {\includegraphics[width=\linewidth]{\@photo}};%
        \end{tikzpicture}%
    \end{minipage}%
    \fi%
    \ifx\@profile\@undefined\else{\par\vspace{\beforeprofileskip}\small\@profile}\fi%
    \endgroup%
    \par\vspace{\titleskip}%
}
%</class>
% \end{comment}
% \end{macro}
%
%
%
% \section{Secciones}
%
% Ahora vamos a trabajar sobre las secciones de tu hoja de vida. Todas las demás partes de la hoja de vida se dividirán en secciones, por ejemplo: «Experiencia Laboral», «Educación», «Premios», etc. Y eventos.
%
% \begin{macro}{\section}
% Es relativamente sencillo este proceso, cada sección nueva de tu hoja de vida se escribe con el comando \cmd\section\marg{título}, con la excepción de que este comando no incluye un argumento opcional (pues en una hoja de vida no hay índice).
% \begin{comment}
%<*class>
\newlength{\beforesectionskip}
\setlength{\beforesectionskip}{1.5\baselineskip\@plus1\baselineskip\@minus1\baselineskip}
\newlength{\sectiontitleskip}
\setlength{\sectiontitleskip}{1\baselineskip\@plus0.5\baselineskip\@minus0.2\baselineskip}
\newlength{\sectionrulewidth}
\setlength{\sectionrulewidth}{2pt}
\newif\if@firstsection
\@firstsectiontrue
\RenewDocumentCommand{\section}{ m }{%
    \@firsteventtrue%
    \@firstskilltrue%
    \if@firstsection\@firstsectionfalse\else\par\vspace{\beforesectionskip}\fi%
    \refstepcounter{section}%
    \currentpdfbookmark{#1}{\thesection}%
    {\LARGE\bfseries\strut\uppercase{#1}}\par\vspace{-0.3\baselineskip}%
    {\hrule height\sectionrulewidth\relax}\par\vspace{\sectiontitleskip}%
}
%</class>
% \end{comment}
% \end{macro}
%
%
%
% \begin{macro}{\event}
% Luego de esto, el siguiente paso sería colocar todos y cada uno de los eventos de esa sección, en caso de que estemos hablando de algo como tu «Experiencia Laboral».
% Para esto simplemente utiliza el comando \cmd\event\marg{nombre}\marg{organización}\marg{ubicación}\marg{fecha}, el cual genera un título con la información que has suministrado.
%
% Luego de esto, si quieres incluir viñetas o \emph{bullet points} mostrando tus logros o actos memorables durante este evento, simplemente utiliza el entorno \texttt{itemize} genérico de {\lmr\LaTeX}.
% Este entorno fue modificado para ser más compacto, y mantener el mismo aspecto visual del resto del contenido.
% \begin{comment}
%<*class>
\newlength{\beforeeventskip}
\setlength{\beforeeventskip}{1\baselineskip\@plus0.5\baselineskip\@minus0.5\baselineskip}
\newlength{\eventtitleskip}
\setlength{\eventtitleskip}{0.5\baselineskip}
\newif\if@firstevent
\NewDocumentCommand{\event}{ m m m m }{%
    \if@firstevent\@firsteventfalse\else\divider[\beforeeventskip]\fi%
    {\large\bfseries#1}\par%
    \IfBlankF{#2}{\textcolor{gray}{\bfseries#2}}\par%
    \IfBlankF{#4}{\parbox{0.5\linewidth}{\color{gray}\tableroutlined{\symbol{"EA52}}\space{#4}}}%
    \IfBlankF{#3}{\parbox{0.5\linewidth}{\color{gray}\tableroutlined{\symbol{"EAE8}}\space{#3}}}%
    \par\vspace{\eventtitleskip}%
}
%</class>
% \end{comment}
% \end{macro}
%
%
%
% \begin{macro}{\skill}
% En caso de que lo que necesites sea registrar tu proficiencia en alguna habilidad o idioma, el comando \cmd\skill\marg{nombre}\marg{cantidad}\oarg{comentario} te puede servir.
% Este comando coloca \meta{nombre} en lo más a la izquierda posible y 5 círculos grises lo más a la derecha posible.
% Esos círculos se cubrirán por otros \meta{cantidad} círculos totalmente negros; este número deber ser un número decimal, si queremos 4 círculos debemos escribir \texttt{4.0} u obtendremos un error.
% El único argumento opcional de este comando nos colocará \meta{comentario} justo debajo de los círculos.
% \begin{comment}
%<*class>
\newlength{\skillheight}
\setlength{\skillheight}{1em}
\newlength{\skillsep}
\setlength{\skillsep}{0.2em}
\newlength{\beforeskillskip}
\setlength{\beforeskillskip}{0.5\baselineskip}
\newif\if@firstskill
\NewDocumentCommand{\skill}{ m m o }{%
    \if@firstskill\@firstskillfalse\else\par\vspace{\skillskip}\fi%
    #1%
    \hfill%
    \begin{tikzpicture}
        \useasboundingbox (-0.5\skillheight,-0.5\skillheight) rectangle (4.5*\skillheight+4*\skillsep,0.5\skillheight);
        \foreach \x in {0,...,4} {
            \fill[black!10] (\skillheight*\x+\skillsep*\x,0) circle (0.5\skillheight);
        }
        \clip (-0.5\skillheight,-0.5\skillheight) rectangle
            (#2*\skillheight+floor{#2}*\skillsep-0.5\skillheight,0.5\skillheight);
        \foreach \x in {0,...,4} {
            \fill (\skillheight*\x+\skillsep*\x,0) circle (0.5\skillheight);
        }
    \end{tikzpicture}%
    \IfValueT{#3}{\par\hfill#3}%
}
%</class>
% \end{comment}
% \end{macro}
%
%
%
% \part{Herramientas adicionales}
%
%
%
% \section{Iconos \label{sec:icons}}
%
% \begin{macro}{\tabler}
% Si quieres incluir un icono, puedes acceder a la fuente de iconos de una manera muy sencilla mediante el comando \cmd\tabler\marg{clase}\marg{icono}.
% En esta clase, por defecto se utilizan los iconos de \href{https://tabler.io/icons}{Tabler}, que posee una colección de más de 6000 iconos distintos para que utilices en cualquier parte de tu documento.
% Estos fueron
%
% Pero no estás restringido solamente a estos, existen muchos otros iconos, tanto gratuitos como pagos, que puedes utilizar.
% Uno de los paquetes más conocidos y populares es \href{https://fontawesome.com}{\emph{FontAwesome}}, este paquete tiene una versión gratuita y otra de pago.
% Sus iconos son bastante bonitos, pero son una fuente pesada (con mucho espacio negro), por lo que su uso es un tanto complicado.
%
% También puedes incluir tus iconos como imágenes, preferiblemente vectoriales.
% Generalmente, los iconos se pueden descargar en formato SVG vectorial, recuerda que el paquete \texttt{graphicx} no recibe gráficos en este formato.
% Puedes utilizar \href{https://inkscape.org}{Inkscape}, \href{https://affinity.studio}{Affinity} (Windows y MacOS), \href{https://vectorpea.com}{Vectorpea} u otro programa de edición vectorial para transformar el formato a PDF (1.5) o EPS.
%
% Puedes utilizar el que prefieras, pero recuerda cambiar de manera adecuada el código de la clase y los \emph{code points} de todos los iconos.
% \begin{comment}
%<*class>
\NewDocumentCommand{\tablerfilled}{ m }{\if@noicons\else\begingroup\tablerfamily\bfseries#1\endgroup\fi}
\NewDocumentCommand{\tableroutlined}{ m }{\if@noicons\else\begingroup\tablerfamily\upshape#1\endgroup\fi}
\NewDocumentCommand{\tabler}{ m m }{\csname tabler#1\endcsname{#2}}
%</class>
% \end{comment}
% \end{macro}
%
%
%
% \section{Divisor}
%
% \begin{macro}{\divider}
% El comando \cmd\divider\oarg{espacio v.} es la definición de la separación entre diferentes eventos.
% Por defecto, el valor de \meta{espacio v.} está definido como |\medskipamount|, que es un espacio bastante considerable.
% Puedes utilizar este comando en caso de que necesites un separador entre diferentes partes de tu hoja de vida, aunque su uso es más interno.
% \begin{comment}
%<*class>
\NewDocumentCommand{\divider}{ O{\medskipamount} }{%
    \par\vspace{0.5#1}%
    \textcolor{gray}{\hbox to \linewidth{\xleaders\hbox to 0.5em{\hfil\rule[0.5ex]{0.25em}{0.4pt}\hfil}\hfill}}%
    \par\vspace{0.5#1}%
}
%</class>
% \end{comment}
% \end{macro}
%
%
%
% \begin{comment}
%<*class>
\AtBeginDocument{%
    \pagestyle{empty}%
    \raggedright%
}
%</class>
% \end{comment}
